\documentclass[11pt,a4paper]{article}

% ---------- Idioma y codificación ----------
\usepackage[utf8]{inputenc}
\usepackage[T1]{fontenc}
\usepackage[spanish,es-noindentfirst]{babel}

% ---------- Página ----------
\usepackage[margin=2.5cm]{geometry}
\usepackage{parskip}
\setlength{\parskip}{0.5\baselineskip}

% ---------- Tipografía y color ----------
\usepackage{xcolor}
\definecolor{codebg}{HTML}{F5F5F5}
\definecolor{notebg}{HTML}{FFF6DC}
\definecolor{noteborder}{HTML}{E0B33C}
\definecolor{linkblue}{HTML}{1A5FB4}

% ---------- Tablas ----------
\usepackage{array}
\usepackage{longtable}
\usepackage{booktabs}

% ---------- Listados de código ----------
\usepackage{listings}
\lstset{
  basicstyle=\ttfamily\small,
  breaklines=true,
  backgroundcolor=\color{codebg},
  frame=single,
  framerule=0pt,
  columns=fullflexible,
  showstringspaces=false,
  xleftmargin=8pt,
  xrightmargin=8pt,
  aboveskip=8pt,
  belowskip=8pt
}

% ---------- Listas ----------
\usepackage{enumitem}
\setlist{itemsep=2pt, topsep=4pt}

% ---------- Encabezados y enlaces ----------
\usepackage{fancyhdr}
\usepackage[hidelinks,colorlinks=true,linkcolor=linkblue,urlcolor=linkblue,citecolor=linkblue]{hyperref}

\pagestyle{fancy}
\fancyhf{}
\fancyhead[L]{Manual de usuario --- Curso de Modelos Estocásticos}
\fancyhead[R]{\thepage}
\renewcommand{\headrulewidth}{0.4pt}

% ---------- Entorno "nota" ----------
\usepackage[most]{tcolorbox}
\newtcolorbox{nota}{
  colback=notebg,
  colframe=noteborder,
  boxrule=1pt,
  arc=3pt,
  left=8pt, right=8pt, top=6pt, bottom=6pt,
  before skip=10pt, after skip=10pt
}

\newcommand{\code}[1]{\texttt{#1}}

\title{\textbf{Manual de usuario}\\[4pt]\large Curso de Modelos Estocásticos --- Plataforma web}
\author{Universidad Nacional de Colombia}
\date{\today}

\begin{document}

\maketitle
\thispagestyle{fancy}

\tableofcontents
\newpage

Esta guía explica cómo usar la plataforma según el tipo de cuenta: \textbf{estudiante}, \textbf{profesor} o \textbf{administrador}.

\section{Acceso a la plataforma}

\begin{enumerate}
  \item Ingresa a la página principal del curso.
  \item Si ya tienes cuenta, haz clic en \textbf{``Iniciar sesión''} (esquina superior derecha) e ingresa tu correo y contraseña.
  \item Si no tienes cuenta, un administrador debe crearla (individualmente o mediante carga masiva por CSV). No hay registro público abierto habilitado en la interfaz actual, aunque existe la ruta de auto-registro a nivel técnico.
  \item Tras iniciar sesión, serás redirigido a la página principal, donde verás tu nombre y un botón de acceso al \textbf{Perfil}.
\end{enumerate}

\subsection*{Cerrar sesión}

Desde la página principal o desde tu Perfil, haz clic en el botón de salida (ícono de flecha/``Salir''). Esto cierra tu sesión en todos los dispositivos donde uses ese mismo token.

\section{Navegación por el contenido del curso}

\begin{itemize}
  \item El \textbf{menú lateral} (ícono $\equiv$) muestra la estructura del curso organizada en \textbf{Partes $\rightarrow$ Secciones $\rightarrow$ Subtemas}.
  \item Haz clic en una Parte para desplegar sus secciones, y en una sección con subtemas para ver sus subsecciones.
  \item Solo verás los temas que el administrador haya marcado como visibles para tu grupo (Grupo 1 o Grupo 2). Si un tema no aparece, puede que aún no esté habilitado para tu grupo.
  \item Dentro de cada tema encontrarás, según lo que el administrador haya publicado:
  \begin{itemize}
    \item \textbf{Texto explicativo.}
    \item \textbf{Video incrustado.}
    \item \textbf{Material descargable} (botón ``Descargar Material'').
    \item \textbf{Tareas y enlaces relacionados.}
  \end{itemize}
  \item El botón de \textbf{modo claro/oscuro} está en la parte superior del menú lateral.
\end{itemize}

\section{Perfil de usuario}

En \textbf{Perfil} puedes ver tu nombre, correo, rol (Estudiante/Profesor/Administrador) y, si aplica, tu grupo asignado.

Los estudiantes y profesores no tienen opciones adicionales en esta pantalla salvo el cierre de sesión y el cambio de tema claro/oscuro.

\section{Panel de administración (solo administradores)}

Al entrar a \textbf{Perfil} con una cuenta de administrador, aparece la sección \textbf{``Panel de Administración''} con tres pestañas.

\subsection{Pestaña ``Cargar CSV''}

Permite crear estudiantes en bloque a partir de un archivo CSV.

\paragraph{Formato requerido del archivo} (columnas, en este orden de nombres):

\begin{lstlisting}[language=]
No, Apellidos, Nombre, Documento, Correo, Grupo
\end{lstlisting}

\begin{itemize}
  \item \code{Grupo} debe ser \code{1} o \code{2}.
  \item La \textbf{contraseña inicial} de cada estudiante creado será su número de \textbf{Documento}. Comunícales esto para que puedan iniciar sesión (actualmente la plataforma no ofrece un flujo de ``olvidé mi contraseña'' ni de cambio de contraseña).
  \item Filas con correo o documento vacío, grupo inválido, o correo ya registrado se omiten; al finalizar la carga se muestra cuántos usuarios se crearon y, si hubo filas omitidas, un aviso con el detalle.
\end{itemize}

\paragraph{Pasos}
\begin{enumerate}
  \item Selecciona el archivo CSV con el botón de carga.
  \item Haz clic en \textbf{``Procesar CSV''}.
  \item Revisa el mensaje de resultado (creados / omitidos / errores).
\end{enumerate}

\subsection{Pestaña ``Usuarios''}

Lista todos los usuarios registrados.

\begin{itemize}
  \item Haz clic en \textbf{``Actualizar''} para refrescar la lista.
  \item Para eliminar un usuario, haz clic en \textbf{``Eliminar''} junto a su nombre y confirma. Esta acción es irreversible y borra también sus sesiones activas.
  \item No es posible eliminar tu propia cuenta desde aquí.
\end{itemize}

\subsection{Pestaña ``Visibilidad''}

Controla qué temas ve cada grupo de estudiantes y permite crear temas nuevos.

\paragraph{Ver/editar visibilidad}
\begin{itemize}
  \item Los temas se agrupan por ``Parte''; haz clic en el nombre de una parte para desplegarla.
  \item Cada tema (y sus subtemas, indentados) muestra dos botones: \textbf{G1} y \textbf{G2}.
  \begin{itemize}
    \item Azul/activo = visible para ese grupo. Gris/inactivo = oculto para ese grupo.
    \item Haz clic en el botón para alternar la visibilidad del tema para ese grupo específico.
  \end{itemize}
\end{itemize}

\paragraph{Sincronizar}
El botón \textbf{``Sincronizar''} actualiza la lista de temas contra el temario oficial del curso: crea los temas que falten y actualiza el título/orden de los existentes. Útil tras cambios en el temario hechos por el equipo de desarrollo. No borra temas existentes.

\paragraph{Crear un tema nuevo}
\begin{enumerate}
  \item Haz clic en \textbf{``+ Crear Tema''}.
  \item Escribe el \textbf{título}; el campo de \textbf{slug} (identificador para la URL) se genera automáticamente, pero puedes editarlo.
  \item Opcionalmente, selecciona una \textbf{Parte} existente o escribe el nombre de una nueva.
  \item Opcionalmente, selecciona un \textbf{tema padre} si este nuevo tema debe aparecer como subtema de otro.
  \item Haz clic en \textbf{``Guardar Tema''}.
\end{enumerate}

\paragraph{Editar el contenido de un tema}
\begin{itemize}
  \item Junto a cada tema en la lista hay un botón \textbf{``Contenido''} que abre el editor de ese tema (\code{/perfil/tema/[id]}).
  \item Desde ahí puedes:
  \begin{itemize}
    \item Cambiar el \textbf{título} del tema (el slug/URL no es editable desde la interfaz, para no romper enlaces existentes).
    \item Añadir bloques de contenido con los botones \textbf{``+ Texto''}, \textbf{``+ Video''}, \textbf{``+ PDF/Material''}:
    \begin{itemize}
      \item \emph{Texto}: cuadro de texto libre. Admite HTML básico y \textbf{fórmulas LaTeX}: \code{\$formula\$} para una fórmula en línea (ej. \code{\$x^2\$}) o \code{\$\$formula\$\$} para una fórmula en bloque, centrada (ej. \code{\$\$E[X] = \textbackslash sum x\_i p\_i\$\$}).
      \item \emph{Video}: URL de \textbf{embed} reproducible en un iframe. Funciona con YouTube (formato \code{https://www.youtube.com/embed/...}) o con Google Drive (formato \code{https://drive.google.com/file/d/ID\_DEL\_ARCHIVO/preview}, no el enlace normal ``Ver''/\code{view}). Si usas Drive, el archivo debe compartirse como ``Cualquier persona con el enlace'' --- si requiere iniciar sesión con una cuenta específica, el video no se reproducirá dentro de la página.
      \item \emph{Material}: URL del archivo PDF u otro material (por ejemplo, \code{/documents/archivo.pdf}).
    \end{itemize}
    \item Eliminar un bloque de contenido con el ícono de papelera.
    \item Añadir/eliminar \textbf{enlaces y tareas} (título + URL) en la sección correspondiente.
    \item Al terminar, haz clic en \textbf{``Guardar Cambios''}.
  \end{itemize}
\end{itemize}

\begin{nota}
Al guardar, la plataforma \textbf{reemplaza por completo} los bloques de contenido y enlaces existentes del tema por los que aparecen en pantalla en ese momento --- revisa que todos los bloques que quieras conservar sigan visibles antes de guardar.
\end{nota}

\section{Preguntas frecuentes}

\begin{description}
  \item[¿Por qué no veo un tema que un compañero de otro grupo sí ve?] Los temas pueden habilitarse de forma independiente para Grupo 1 y Grupo 2. Consulta con el administrador si crees que debería estar visible para tu grupo.
  \item[Olvidé mi contraseña, ¿cómo la recupero?] Actualmente no existe un flujo de autoservicio para restablecer contraseña. Contacta al administrador del curso para que la restablezca manualmente (creando de nuevo el usuario o mediante acceso directo a la base de datos).
  \item[Subí el CSV y algunos estudiantes no se crearon, ¿por qué?] Revisa el mensaje de errores que se muestra tras la carga: las causas más comunes son correo o documento vacíos, un valor de grupo distinto de 1 o 2, o un correo que ya existe en el sistema.
  \item[¿Puedo deshacer la eliminación de un usuario o de contenido de un tema?] No. Tanto la eliminación de usuarios como el guardado de cambios en el editor de temas son operaciones definitivas; no hay papelera ni historial de versiones.
\end{description}

\end{document}
